\subsection{The Square Root Method for Quadratics}
\begin{definition}[Square Root Method]
	A technique used to solve quadratic equations by isolating the squared variable term, then taking the square root of both sides.
\end{definition}
\begin{example}[Square Root Method]
	\hfill

	
	\textbf{All of these has 2 solutions.}
	\begin{enumerate}
		\item $f(x) = x^2  - 18$ \\
			$x^2 - 18 = 0$ \\
			$x^2 = 18$ \\
			$\sqrt{x^2} = \pm \sqrt{18}$ \\
			$x = \pm 3\sqrt{2}$ \\
			$x = 3\sqrt{2} \land x = -3\sqrt{2}$
		\item $f(x) = 3x^2 - 33$ \\
			$3x^2 - 33 = 0$ \\
			$3x^2 = 33$ \\
			$x^2 = 11$ \\
			$\sqrt{x^2} = \pm \sqrt{11}$ \\
			$x = \pm \sqrt{11}$ \\
			$x = \sqrt{11} \land x = -\sqrt{11}$
		\item $g(x) = (x + 2)^2 - 1$ \\
			$(x + 2)^2 - 1 = 0$ \\
			$(x + 2)^2 = 1$ \\
			$\sqrt{(x + 2)^2} = \pm \sqrt{1}$ \\
			$x + 2 = \pm 1$ \\
			$x = -2 \pm 1$ \\
			$x = -2 - 1 = -1 \land x = -2 + 1 = -3$
		\item $h(x) = (2x + 3)^2 - 32$ \\
			$(2x + 3)^2 - 32 = 0$ \\
			$(2x + 3)^2 = 32$ \\
			$\sqrt{(2x + 3)^2} = \pm \sqrt{32}$ \\
			$2x + 3 = 4\sqrt{2}$ \\
			$2x + 3 = \pm 4\sqrt{2}$ \\
			$2x = -3 \pm 4\sqrt{2}$ \\
			$x = \frac{-3 \pm 4\sqrt{2}}{2}$ \\
			$x = \frac{-3 - 4\sqrt{2}}{2} \land x = \frac{-3 + 4\sqrt{2}}{2}$
	\end{enumerate}
\end{example}

\begin{definition}[Imaginary Numbers]
	An imaginary number is any number that can be written as a real number multiplied by the imaginary unit $i$ where $i = \sqrt{-1}$.
\end{definition}

\begin{example}[Special Cases]
	\hfill
	\begin{enumerate}
		\item \textbf{1 Solution.} \\
		 $f(x) = (x - 7)^2$ \\
			$(x - 7)^2 = 0$ \\
			$\sqrt{(x - 7)^2} = \pm \sqrt{0}$ \\
			$x - 7 = \pm 0$ \\
			$x = 7 \pm 0$ \\
			$x = 7 - 0 = 7 \land x = 7 + 0 = 7$
		\item \textbf{No solution (Imaginary Number). }\\
		 $f(x) = (3x - 2)^2 + 75$ \\
			$(3x - 2)^2 + 75 = 0$ \\
			$(3x - 2)^2 = -75$ \\
			$\sqrt{(3x - 2)^2} = \pm \sqrt{-75}$ \\
			$3x - 2 = \pm 5i\sqrt{3}$ \\
			$x = \frac{2 \pm 5i\sqrt{3}}{3}$ \\
			$x = \frac{2 + 5i\sqrt{3}}{3} \land x = \frac{2 + 5i\sqrt{3}}{3}$ \\
			Can also be written as \\
			$x = \frac{2}{3} + \frac{5\sqrt{3}}{3}i \land x = \frac{2}{3} - \frac{5\sqrt{3}}{3}i$
	\end{enumerate}
\end{example}
